\documentclass[sigconf]{acmart}
%% ACM style requires these packages
\usepackage{graphicx}
\usepackage{booktabs} % for professional tables
%% ACM metadata
\acmConference[Conference'26]{ACM Conference}{June 2026}{City, State, Country}
\acmYear{2026}
\copyrightyear{2026}
\acmDOI{XXXXXXX.XXXXXXX}
\acmISBN{978-1-4503-XXXX-X/26/06}
%% Rights management (choose one)
%\setcopyright{acmcopyright}
%\setcopyright{acmlicensed}
\setcopyright{rightsretained}
%%
%% The majority of ACM publications use the "sigconf" format.
%% Other options: sigplan, acmtog, acmsmall, acmlarge
%%
\begin{document}
%%
%% Title and authors
%%
\title{Using Supervised and Unsupervised Machine Learning Methods to Classify XRF Samples}
\author{Larry Griffith}
\affiliation{%
  \institution{Northern Arizona University}
  \city{Flagstaff}
  \state{Arizona}
  \country{United States}
}
\email{ldg246@nau.edu}
\author{Dr. Kayeleigh Sharp}
\affiliation{%
  \institution{Northern Arizona University}
  \city{Flagstaff}
  \state{Arizona}
  \country{United States}
}
\email{kayeleigh.sharp@nau.edu}
%%
%% Abstract (required by ACM)
%%
\begin{abstract}
% SCAFFOLD - write this LAST, ~150-200 words. One sentence per move:
%   1. Context: portable XRF (pXRF) produces elemental readings on archaeological
%      samples, but assigning each reading to a material class (pottery, soil,
%      slag, metal, etc.) is manual and subjective.
%   2. Gap: hand-labeling is slow and inconsistent; labeled data is scarce.
%   3. What you did: you built ArchaeoSight, a desktop tool that classifies pXRF
%      samples using BOTH a supervised path (gradient-boosted decision trees) and
%      an unsupervised path (autoencoder + HDBSCAN), plus geostatistical kriging
%      and a "next dig" recommender, and a Monte Carlo synthetic-data generator.
%   4. Key numbers: report headline multi-class and binary (soil/non-soil)
%      accuracy for the GBDT, silhouette / cluster count for the unsupervised path.
%   5. Why it matters: faster, reproducible material classification and
%      data-driven excavation targeting.
\end{abstract}
%%
%% Keywords (required by ACM)
%%
\keywords{machine learning, XRF, classification, archaeology, supervised learning,
  unsupervised learning, gradient boosting, autoencoder, HDBSCAN, kriging,
  geostatistics, active sampling, synthetic data}
%%
%% CCS Concepts (required by ACM — generate at https://dl.acm.org/ccs)
%%
% SCAFFOLD - the one below covers ML. Consider ALSO adding:
%   Applied computing~Arts and humanities  (archaeology angle)
%   Computing methodologies~Cluster analysis  (HDBSCAN path)
\begin{CCSXML}
<ccs2012>
  <concept>
    <concept_id>10010147.10010178</concept_id>
    <concept_desc>Computing methodologies~Machine learning</concept_desc>
    <concept_significance>500</concept_significance>
  </concept>
</ccs2012>
\end{CCSXML}
\ccsdesc[500]{Computing methodologies~Machine learning}
\maketitle
%%
%% Body
%%
\section{Introduction}
% SCAFFOLD - the archaeology-first framing. Cover, in order:
%   - What pXRF data is: each sample is a vector of element concentrations
%     (the dataset uses Ag, Al, As, Au, Ca, Cr, Cu, Fe, K, Mg, Mn, Ni, P, Pb,
%     Sr, Ti, V, Zn) plus spatial coordinates (X_Coord, Y_Coord) and context
%     metadata (BAG, CNTXT, AreaFunction, level).
%   - The problem: archaeologists must decide what MATERIAL each reading is
%     (the 'material' column: pottery, soil, etc.). Doing this by eye is slow,
%     needs expertise, and is inconsistent between analysts.
%   - The motivation: a tool that classifies automatically AND helps decide
%     where to dig next would speed up fieldwork and reduce subjectivity.
%   - State the research questions explicitly, e.g.:
%       RQ1: Can supervised ML (GBDT) classify pXRF samples by material with
%            usable accuracy on this dataset?
%       RQ2: Can unsupervised ML (autoencoder + HDBSCAN) recover meaningful
%            material groupings WITHOUT labels?
%       RQ3: Can geostatistics (kriging) + an acquisition score recommend
%            high-value next excavation locations?
%   - End with a one-paragraph summary of contributions and a roadmap sentence.

\subsection{Methods}
% SCAFFOLD - HIGH-LEVEL only here (the detailed version lives in Section 3).
% Name the four pipelines the tool implements and one line on each:
%   1. Supervised classification - gradient-boosted decision trees on raw
%      element concentrations (GradientBoostedDecisionTreePage.py).
%   2. Unsupervised classification - autoencoder -> PCA -> HDBSCAN clustering,
%      then clusters mapped to materials (ClusteringPage.py).
%   3. Spatial interpolation - ordinary/universal kriging of an element's
%      concentration across the site (KrigingPage.py).
%   4. Next-dig recommendation - kriging value + uncertainty combined into a
%      priority score to rank candidate excavation cells (NextDigPage.py).
%   Mention the Monte Carlo synthetic-sample generator as supporting data
%   augmentation (montecarlo_samples.py).

\section{Background}
% SCAFFOLD - what led here + prior work. Suggested subtopics:
%   - pXRF in archaeology: how it works, why elemental fingerprints relate to
%     material/provenance, known limitations (matrix effects, surface vs bulk).
%   - Prior ML-on-XRF work: cite studies that classify materials or sourcing
%     from XRF/pXRF using ML (find 3-5 references for references.bib).
%   - Gradient boosting background: brief, cite the GBM / boosting literature.
%   - Representation learning + density clustering: autoencoders for
%     dimensionality reduction; HDBSCAN as density-based, noise-aware,
%     no-fixed-k clustering (cite the HDBSCAN paper).
%   - Geostatistics: kriging and the variogram, why it suits spatially
%     autocorrelated concentration data (cite a geostatistics reference).
%   - Active/sequential sampling: explore-exploit acquisition for choosing
%     informative new measurement locations (cite Bayesian optimization /
%     active learning).
%   - Synthetic data / copulas: why class-conditional augmentation helps with
%     small, imbalanced labeled sets.

\section{Methods}

\subsection{Implementation}
% SCAFFOLD - describe the system and EACH model precisely. The tool is a PyQt6
% desktop app; each method is a page/worker. Document these as they are coded:
%
% (a) Supervised - Gradient-Boosted Decision Trees (GradientBoostedDecisionTreePage.py)
%   - Features: only true element columns are auto-selected (intersection of
%     dataframe columns with the periodic-table element list).
%   - Preprocessing: NaN -> 0, negatives clipped to 0; labels via LabelEncoder;
%     features standardized with StandardScaler.
%   - Model: sklearn GradientBoostingClassifier with user-set hyperparameters
%     (n_estimators, learning_rate, max_depth).
%   - Validation: 5-fold StratifiedKFold cross-validation (report mean +- 2 std),
%     then an 80/20 stratified train/test split for the final model.
%   - Evaluation: multi-class accuracy, plus a BINARY soil/non-soil accuracy
%     derived by thresholding the predicted soil probability at 0.5; confusion
%     matrix; per-class precision/recall/F1; feature importances per element.
%   - Export: pickle or ONNX (skl2onnx) for portable inference.
%
% (b) Unsupervised - Autoencoder + HDBSCAN (ClusteringPage.py)
%   - Preprocessing: numeric coercion, NaN -> 0, clip negatives, StandardScaler.
%     If labels exist, stratified split on the binary soil/non-soil label.
%   - Representation: a Keras/TensorFlow autoencoder compresses the element
%     vector to a 'latent_dim' bottleneck; Adam(1e-3), MSE loss, EarlyStopping
%     on val_loss (patience 10, restore best). Save the encoder + training-loss
%     curve.
%   - Reduce: PCA (<=5 components) on the latent vectors.
%   - Cluster: HDBSCAN (min_cluster_size) on PCA(latent); label -1 = noise.
%     approximate_predict assigns the held-out/test points; weighted cluster
%     centroids stored for a KNN-style approximation.
%   - Label mapping (when ground truth is present): majority vote maps each
%     cluster to its dominant material; train RandomForest classifiers on the
%     latent space (binary soil/non-soil and multi-class) for evaluation.
%   - Evaluation: silhouette score, number of clusters, cluster size counts,
%     PCA(latent) scatter colored by cluster.
%   - Export: encoder + HDBSCAN bundle, ONNX via tf2onnx / skl2onnx.
%
% (c) Spatial interpolation - Kriging (KrigingPage.py)
%   - Inputs: X_Coord, Y_Coord, and one chosen element as Z.
%   - Optional ln(1+Z) log-transform; kriging math runs in transformed space,
%     interpolation map back-transformed with expm1, variance/variogram stay in
%     log space.
%   - Build an nx-by-ny interpolation grid with margins.
%   - pykrige OrdinaryKriging or UniversalKriging; variogram models available:
%     linear, power, gaussian, spherical, exponential, hole-effect; nlags and
%     manual vs auto-fit variogram parameters configurable.
%   - Outputs: interpolated concentration surface, kriging variance surface,
%     and the fitted variogram (model parameters extracted).
%
% (d) Next-dig recommender - active sampling (NextDigPage.py)
%   - Run an ordinary-kriging surrogate to get predicted value z_grid and
%     kriging std-dev sigma = sqrt(variance) on the grid.
%   - Normalize value and sigma; priority = (1 - w)*value_norm + w*sigma_norm,
%     where w is the explore/exploit weight (0 = exploit high values, 1 =
%     explore high uncertainty).
%   - Mask to the surveyed area (convex hull / Delaunay) and exclude an "avoid
%     radius" around existing samples (cKDTree) so picks are not resamples.
%   - Rank remaining candidate cells and emit the top-k as ranked dig
%     recommendations with predicted value, sigma, and score; render a priority
%     heatmap and a site-image overlay.
%
% (e) Synthetic data - Monte Carlo per-class generator (montecarlo_samples.py)
%   - Per material class, fit a Gaussian copula: each feature's marginal is
%     preserved exactly via inverse empirical CDF sampling (no values outside
%     the observed range); joint structure from a PSD-corrected correlation
%     matrix.
%   - Non-numeric / carry-through columns (e.g. coordinates, IDs) are
%     bootstrap-resampled within the class to keep schema and categoricals.
%   - Controls: --n-per-class or --multiplier, --exclude, --include-original;
%     report correlation drift between original and synthetic as a quality check.

\subsection{Datasets}
% SCAFFOLD - describe the actual data files:
%   - AllPXRF_FINAL_14Oct.csv: the real pXRF readings. State row count, the
%     element columns measured (Ag, Al, As, Au, Ca, Cr, Cu, Fe, K, Mg, Mn, Ni,
%     P, Pb, Sr, Ti, V, Zn), the 'material' label column and its class
%     distribution, spatial columns (X_Coord, Y_Coord), and context fields
%     (BAG, CNTXT, AreaFunction, level, pXRFno).
%   - AllPXRF_FINAL_14Oct_synthetic.csv and combined.csv: Monte Carlo augmented
%     / merged sets - explain how/why they were generated and whether models
%     were trained on real, synthetic, or combined data.
%   - Note preprocessing common to all models (negative values -> 0, NaN -> 0).
%   - Report class imbalance and how it was handled (stratification, synthetic
%     augmentation). A small table of class counts would go well here.

\subsection{Successes and Challenges}
% SCAFFOLD - be candid and concrete:
%   Successes: which method classified best and the numbers; whether
%   unsupervised clusters lined up with real materials; useful element
%   importances (which elements drive class separation); whether kriging
%   surfaces matched archaeological expectation; sensible next-dig picks.
%   Challenges: class imbalance; noise / -1 HDBSCAN labels; autoencoder
%   sensitivity to latent_dim and epochs; variogram fitting failures on sparse
%   spatial data; risk that synthetic data leaks distribution assumptions;
%   matrix effects in pXRF; reproducibility (fixed seed = 42 noted in code).

\section{Discussion}
% SCAFFOLD - interpret, do not just restate results. Tie back to the RQs:
%   - Supervised vs unsupervised trade-off: GBDT is accurate but needs labels;
%     the autoencoder+HDBSCAN path finds structure without them - what does that
%     buy an archaeologist with little labeled data?
%   - What do the cluster-to-material mappings and feature importances tell us
%     archaeologically (which elements signal which materials)?
%   - How the kriging + next-dig pipeline changes fieldwork decisions
%     (exploit known hotspots vs explore uncertain areas).
%   - Threats to validity: site-specificity, transfer to other assemblages,
%     dependence on label quality, synthetic-data realism.

\section{Future Work}
% SCAFFOLD - directions:
%   - Try other supervised models (XGBoost/LightGBM, calibrated probabilities).
%   - Semi-supervised: use HDBSCAN clusters to propose labels for active
%     re-labeling.
%   - Close the loop: feed next-dig recommendations back into model retraining
%     (true active learning).
%   - Uncertainty-aware classification and conformal prediction.
%   - Validate synthetic data with downstream-accuracy and statistical tests.
%   - Cross-site generalization and field deployment of the desktop tool.

\section{Conclusion}
% SCAFFOLD - 1 short paragraph. Restate the contribution (a desktop tool pairing
% supervised GBDT and unsupervised autoencoder+HDBSCAN classification of pXRF
% samples, with kriging-based spatial interpolation and a data-driven next-dig
% recommender), the headline result(s), and the takeaway for archaeological
% practice. No new information.
%%
%% Acknowledgments
%%
\begin{acks}
  % TODO: Add acknowledgments, funding sources, etc.
  % SCAFFOLD - advisor (Dr. Sharp), NAU, the excavation/dataset providers,
  % any grant numbers.
\end{acks}
%%
%% Bibliography
%% Use BibTeX with the ACM reference format
%%
% SCAFFOLD - create references.bib with entries for: pXRF in archaeology,
% gradient boosting, autoencoders, HDBSCAN, kriging/geostatistics, active
% learning / Bayesian optimization, Gaussian copulas / synthetic data, and the
% core libraries used (scikit-learn, TensorFlow/Keras, PyKrige, hdbscan).
\bibliographystyle{ACM-Reference-Format}
\bibliography{references} % expects a references.bib file
\end{document}
